\subsection*{CU-39: Comprar y abrir una caja}
\addcontentsline{toc}{subsection}{CU-39: Comprar y abrir una caja}
\label{cu-39}

\begin{fichacu}{CU-39}{Comprar y abrir una caja}
  \crudFila{Responsable}{Ángel Gabriel Aguilar Hernández}
  \crudFila{Actor principal}{Jugador}
  \crudFila{Actores de sistema}{Servidor de partidas, Base de datos}
  \crudFila{Prototipos}{14a, 14b, 14c, 14d, 14e, 18c}
  \crudFila{Objetivo}{Adquirir una caja con monedas y obtener sus tres objetos cosméticos, con las probabilidades a la vista, sin repetidos imposibles de convertir y con la garantía de rareza cumplida.}
  \crudFila{Disparador}{El Jugador selecciona una caja en la \texttt{GUI\_\allowbreak{}Shop}, o abre una caja ganada jugando desde su inventario.}
  \textbf{Precondiciones} & \cuCelda{\begin{cureglas}
      \item \textbf{\mbox{PRE-01}} El Jugador tiene una \textit{Sesion} abierta y su token es válido.
      \item \textbf{\mbox{PRE-02}} El catálogo contiene al menos un \textit{TipoCaja} activo con su tabla de probabilidades publicada.
      \item \textbf{\mbox{PRE-03}} El Servidor de partidas está en ejecución y tiene conexión disponible con la Base de datos.
    \end{cureglas}} \\ \hline
  \textbf{Flujo normal} & \cuCelda{\begin{cuenum}[start=1]
      \item El SJM despliega la sección de cajas de la \texttt{GUI\_\allowbreak{}Shop} con los tres tipos, su precio y el acceso a su contenido y probabilidades.
      \item El Jugador consulta el contenido y las probabilidades de un tipo (\mbox{RN-01}). \textit{(Ver \mbox{FA-01})}
      \item El Jugador da click en ``Comprar caja''.
      \item El SJM despliega la \texttt{GUI\_\allowbreak{}BoxPurchaseConfirm} con el precio, el saldo resultante y, de forma visible, cuántas cajas lleva sin objeto raro y cuándo se activa la garantía (\mbox{RN-05}). \textit{(Ver \mbox{FA-02})}
      \item El Jugador confirma la compra.
      \item El SJM envía el mensaje \texttt{BUY\_\allowbreak{}BOX\_\allowbreak{}REQUEST} con el tipo de caja, el precio mostrado y el token de sesión. \textit{(Ver \mbox{EX-03}, \mbox{EX-04})}
      \item El Servidor de partidas verifica que el token corresponda a una \textit{Sesion} abierta. \textit{(Ver \mbox{FA-03})}
      \item El Servidor de partidas verifica que el \textit{TipoCaja} exista, esté activo y que el precio coincida con el del catálogo. \textit{(Ver \mbox{FA-04})}
    \end{cuenum}} \\
   & \cuCelda{\begin{cuenum}[start=9]
      \item El Servidor de partidas verifica que el saldo del \textit{Usuario} alcance el precio. \textit{(Ver \mbox{FA-05})}
      \item El Servidor de partidas inicia una transacción sobre la Base de datos.
      \item El Servidor de partidas crea la \textit{Caja} del \textit{Usuario}, con su tipo, el origen \texttt{COMPRA}, la fecha y el estado \texttt{SIN\_\allowbreak{}ABRIR}. \textit{(Ver \mbox{EX-01})}
      \item El Servidor de partidas crea el \textit{MovimientoMoneda} de tipo \texttt{COMPRA\_\allowbreak{}CAJA} con el importe en negativo y descuenta el saldo.
      \item El Servidor de partidas confirma la transacción. \textit{(Ver \mbox{EX-02})}
      \item El Servidor de partidas responde con \texttt{BUY\_\allowbreak{}BOX\_\allowbreak{}OK} y la caja creada.
      \item El SJM ofrece abrirla ahora o dejarla en el inventario. \textit{(Ver \mbox{FA-06})}
      \item El Jugador selecciona ``Abrir''.
    \end{cuenum}} \\
   & \cuCelda{\begin{cuenum}[start=17]
      \item El SJM envía el mensaje \texttt{OPEN\_\allowbreak{}BOX\_\allowbreak{}REQUEST} con el identificador de la caja.
      \item El Servidor de partidas verifica que la \textit{Caja} pertenezca al \textit{Usuario} y esté \texttt{SIN\_\allowbreak{}ABRIR}. \textit{(Ver \mbox{FA-07})}
      \item El Servidor de partidas verifica que la caja no haya caducado, si su origen es de recompensa (\mbox{RN-06}). \textit{(Ver \mbox{FA-08})}
      \item El Servidor de partidas sortea tres objetos del \textit{TipoCaja} conforme a su tabla de probabilidades, descartando los que el \textit{Usuario} ya posea por duplicado y aplicando la garantía de rareza si procede (\mbox{RN-03}, \mbox{RN-04}, \mbox{RN-05}). \textit{(Ver \mbox{EX-05})}
      \item El Servidor de partidas inicia una transacción sobre la Base de datos.
      \item El Servidor de partidas crea un \textit{UsuarioObjeto} por cada objeto sorteado que el \textit{Usuario} no poseía, con el origen \texttt{CAJA}. \textit{(Ver \mbox{EX-01})}
      \item El Servidor de partidas convierte en monedas cada objeto sorteado que el \textit{Usuario} ya poseía, creando un \textit{MovimientoMoneda} de tipo \texttt{CONVERSION} con el importe en positivo y sumándolo al saldo (\mbox{RN-03}).
      \item El Servidor de partidas crea una \textit{CajaContenido} por cada uno de los tres objetos, registrando qué salió y si se convirtió, para que el resultado sea consultable después (\mbox{RN-07}).
    \end{cuenum}} \\
   & \cuCelda{\begin{cuenum}[start=25]
      \item El Servidor de partidas actualiza en el \textit{Usuario} el contador de cajas sin objeto raro, poniéndolo a cero si salió uno e incrementándolo si no (\mbox{RN-05}).
      \item El Servidor de partidas marca la \textit{Caja} como \texttt{ABIERTA} con su fecha de apertura.
      \item El Servidor de partidas confirma la transacción. \textit{(Ver \mbox{EX-02})}
      \item El Servidor de partidas responde con \texttt{OPEN\_\allowbreak{}BOX\_\allowbreak{}OK} y los tres objetos, indicando cuáles se convirtieron y por cuántas monedas.
      \item El SJM despliega la animación de apertura sobre el tablero, con las tres cartas boca abajo. \textit{(Ver \mbox{FA-09})}
      \item El SJM muestra el resultado con lo que salió, lo convertido y el saldo actualizado.
      \item Fin del caso de uso.
    \end{cuenum}} \\ \hline
  \textbf{Flujos alternos} & \cuCelda{\textbf{\mbox{FA-01}: El Jugador consulta el contenido completo}\par
    \begin{cuenum}[start=1]
      \item El SJM despliega la lista completa del \textit{TipoCaja} con cada objeto y su probabilidad exacta.
      \item El flujo continúa en el paso 2 del FN. Las probabilidades se muestran \textbf{antes} de comprar, siempre y sin tener que buscarlas (\mbox{RN-01}).
    \end{cuenum}} \\
   & \cuCelda{\textbf{\mbox{FA-02}: El Jugador cancela la compra}\par
    \begin{cuenum}[start=1]
      \item El Jugador selecciona ``Cancelar''.
      \item El SJM cierra la confirmación sin enviar nada al Servidor de partidas.
      \item El flujo continúa en el paso 1 del FN.
    \end{cuenum}} \\
   & \cuCelda{\textbf{\mbox{FA-03}: La sesión ya no es válida}\par
    \begin{cuenum}[start=1]
      \item El Servidor de partidas responde con \texttt{SESSION\_\allowbreak{}INVALID}.
      \item El SJM muestra la \texttt{GUI\_\allowbreak{}MessageWarning} y despliega la \texttt{GUI\_\allowbreak{}Login}.
      \item Fin del caso de uso.
    \end{cuenum}} \\
   & \cuCelda{\textbf{\mbox{FA-04}: El tipo de caja ya no está disponible o cambió de precio}\par
    \begin{cuenum}[start=1]
      \item El Servidor de partidas responde con \texttt{BUY\_\allowbreak{}BOX\_\allowbreak{}FAILED} y el motivo correspondiente, incluyendo el precio vigente si es el caso.
      \item El Servidor de partidas registra la discrepancia de precio en la bitácora del servidor.
      \item El SJM recarga la sección de cajas.
      \item El flujo continúa en el paso 1 del FN.
    \end{cuenum}} \\
   & \cuCelda{\textbf{\mbox{FA-05}: El saldo no alcanza}\par
    \begin{cuenum}[start=1]
      \item El Servidor de partidas responde con \texttt{BUY\_\allowbreak{}BOX\_\allowbreak{}FAILED} y el motivo \texttt{SALDO\_\allowbreak{}INSUFICIENTE}, con el importe que falta.
      \item El SJM lo indica, sin ofrecer ninguna forma de obtener monedas con dinero real: no existe (\mbox{RN-02}).
      \item El flujo continúa en el paso 1 del FN.
    \end{cuenum}} \\
   & \cuCelda{\textbf{\mbox{FA-06}: El Jugador deja la caja sin abrir}\par
    \begin{cuenum}[start=1]
      \item El SJM cierra el aviso y la \textit{Caja} permanece en el inventario con estado \texttt{SIN\_\allowbreak{}ABRIR}.
      \item Fin del caso de uso. El Jugador puede abrirla más tarde, retomando el flujo en el paso 16.
    \end{cuenum}} \\
   & \cuCelda{\textbf{\mbox{FA-07}: La caja ya fue abierta}\par
    \begin{cuenum}[start=1]
      \item El Servidor de partidas responde con \texttt{OPEN\_\allowbreak{}BOX\_\allowbreak{}FAILED} y el motivo \texttt{CAJA\_\allowbreak{}YA\_\allowbreak{}ABIERTA}, junto con el resultado que tuvo.
      \item El SJM muestra ese resultado en lugar de la animación.
      \item Fin del caso de uso.
    \end{cuenum}} \\
   & \cuCelda{\textbf{\mbox{FA-08}: La caja caducó}\par
    \begin{cuenum}[start=1]
      \item El Servidor de partidas marca la \textit{Caja} como \texttt{CADUCADA}.
      \item El Servidor de partidas responde con \texttt{OPEN\_\allowbreak{}BOX\_\allowbreak{}FAILED} y el motivo \texttt{CAJA\_\allowbreak{}CADUCADA}.
      \item El SJM lo indica y retira la caja del inventario. Las cajas ganadas jugando caducan a los catorce días; las compradas, no (\mbox{RN-06}).
      \item Fin del caso de uso.
    \end{cuenum}} \\
   & \cuCelda{\textbf{\mbox{FA-09}: El Jugador se salta la animación}\par
    \begin{cuenum}[start=1]
      \item El Jugador selecciona ``Saltar'' durante la apertura.
      \item El SJM muestra el resultado de inmediato. El sorteo ya estaba hecho en el servidor: la animación no influye en él (\mbox{RN-08}).
      \item El flujo continúa en el paso 30 del FN.
    \end{cuenum}} \\
   & \cuCelda{\textbf{\mbox{FA-10}: Los tres objetos sorteados ya se poseían}\par
    \begin{cuenum}[start=1]
      \item El Servidor de partidas convierte los tres en monedas conforme a \mbox{RN-03}.
      \item El SJM muestra el resultado con las tres conversiones y el total obtenido, sin ocultarlo ni disfrazarlo.
      \item El flujo continúa en el paso 30 del FN.
    \end{cuenum}} \\ \hline
  \textbf{Excepciones} & \cuCelda{\textbf{\mbox{EX-01}: Fallo al escribir en la base de datos}\par
    \begin{cuenum}[start=1]
      \item El Servidor de partidas revierte la transacción, de modo que no quede una caja cobrada sin crearse, ni abierta sin entregar su contenido (\mbox{RN-09}).
      \item El Servidor de partidas registra la excepción con un identificador de incidencia y responde con \texttt{SERVER\_\allowbreak{}ERROR}.
      \item El SJM muestra la \texttt{GUI\_\allowbreak{}MessageError} con el identificador y recarga el saldo y el inventario.
      \item Fin del caso de uso.
    \end{cuenum}} \\
   & \cuCelda{\textbf{\mbox{EX-02}: Fallo al confirmar la transacción}\par
    \begin{cuenum}[start=1]
      \item El Servidor de partidas trata la transacción como fallida y registra la excepción señalando que el estado es indeterminado.
      \item Si la caja llegó a cobrarse pero no a crearse, el Servidor de partidas devuelve el importe al detectarlo, creando un \textit{MovimientoMoneda} de tipo \texttt{DEVOLUCION} (\mbox{RN-10}).
      \item El SJM muestra la \texttt{GUI\_\allowbreak{}MessageError} indicando que revise su inventario y su saldo.
      \item Fin del caso de uso.
    \end{cuenum}} \\
   & \cuCelda{\textbf{\mbox{EX-03}: Se agota el tiempo de espera de la respuesta}\par
    \begin{cuenum}[start=1]
      \item El SJM no reenvía la compra ni la apertura: podrían aplicarse dos veces.
      \item El SJM solicita el saldo y el inventario actualizados.
      \item Fin del caso de uso.
    \end{cuenum}} \\
   & \cuCelda{\textbf{\mbox{EX-04}: La conexión se interrumpe}\par
    \begin{cuenum}[start=1]
      \item El SJM muestra la barra de conexión perdida.
      \item Al recuperarla, el SJM recarga el inventario, que dirá si la caja existe y si está abierta.
      \item Fin del caso de uso.
    \end{cuenum}} \\
   & \cuCelda{\textbf{\mbox{EX-05}: El sorteo no puede completarse}\par
    \begin{cuenum}[start=1]
      \item El Servidor de partidas detecta que el \textit{TipoCaja} no tiene objetos suficientes o que su tabla de probabilidades no suma lo que debe.
      \item El Servidor de partidas no abre la caja, la deja \texttt{SIN\_\allowbreak{}ABRIR} y registra la incidencia en la bitácora del servidor para su revisión.
      \item El Servidor de partidas responde con \texttt{SERVER\_\allowbreak{}ERROR}.
      \item El SJM indica que la caja no pudo abrirse y que sigue en su inventario.
      \item Fin del caso de uso.
    \end{cuenum}} \\ \hline
  \textbf{Postcondiciones} & \cuCelda{\begin{cureglas}
      \item \textbf{\mbox{POST-01}} Si el caso termina por el FN, existe una \textit{Caja} con estado \texttt{ABIERTA}, su tipo, su origen y sus fechas de obtención y apertura.
      \item \textbf{\mbox{POST-02}} Si el caso termina por el FN, existen tres \textit{CajaContenido} que dicen exactamente qué salió y qué se convirtió.
      \item \textbf{\mbox{POST-03}} Si el caso termina por el FN, el \textit{Usuario} posee, mediante \textit{UsuarioObjeto}, todos los objetos sorteados que no tenía.
      \item \textbf{\mbox{POST-04}} Si el caso termina por el FN, existe un \textit{MovimientoMoneda} de tipo \texttt{COMPRA\_\allowbreak{}CAJA} en negativo y uno de tipo \texttt{CONVERSION} en positivo por cada repetido.
      \item \textbf{\mbox{POST-05}} Si el caso termina por el FN, el saldo del \textit{Usuario} sigue siendo la suma exacta de sus \textit{MovimientoMoneda}.
      \item \textbf{\mbox{POST-06}} Si el caso termina por el FN, el contador de cajas sin objeto raro quedó a cero si salió uno, o incrementado si no.
      \item \textbf{\mbox{POST-07}} En ningún caso el \textit{Usuario} posee dos veces el mismo \textit{ObjetoCosmetico}.
    \end{cureglas}} \\
   & \cuCelda{\begin{cureglas}
      \item \textbf{\mbox{POST-08}} Si el caso termina por el \mbox{FA-06}, la caja existe \texttt{SIN\_\allowbreak{}ABRIR} y el importe ya se cobró.
      \item \textbf{\mbox{POST-09}} Si el caso termina por cualquier excepción, o la operación quedó completa o no quedó nada de ella, y ninguna caja cobrada se perdió sin abrir ni devolver.
    \end{cureglas}} \\ \hline
  \textbf{Reglas de negocio} & \cuCelda{\begin{cureglas}
      \item \textbf{\mbox{RN-01}} Las probabilidades de cada \textit{TipoCaja} se muestran siempre antes de comprar, con el valor exacto de cada objeto.
      \item \textbf{\mbox{RN-02}} Las cajas solo se compran con monedas del juego. No hay dinero real en ningún punto del sistema.
      \item \textbf{\mbox{RN-03}} Un objeto repetido se convierte en monedas de forma automática. El Jugador no puede quedarse con un duplicado inútil.
      \item \textbf{\mbox{RN-04}} Nunca sale un objeto que el \textit{Usuario} ya tenga dos veces, porque la posesión es única: el duplicado se convierte en el mismo acto.
      \item \textbf{\mbox{RN-05}} A las diez cajas sin un objeto raro, la siguiente lo garantiza. El contador vive en la cuenta y se muestra antes de comprar.
      \item \textbf{\mbox{RN-06}} Las cajas ganadas jugando caducan a los catorce días. Las compradas no caducan.
      \item \textbf{\mbox{RN-07}} El contenido de cada caja abierta queda registrado. Sin ese registro, un jugador que reclame lo que le salió no tendría cómo comprobarlo, y el sistema tampoco.
    \end{cureglas}} \\
   & \cuCelda{\begin{cureglas}
      \item \textbf{\mbox{RN-08}} El sorteo lo hace el servidor antes de la animación. La animación representa un resultado ya decidido y saltarla no lo cambia.
      \item \textbf{\mbox{RN-09}} La compra y la apertura son dos transacciones distintas, cada una atómica. Entre ellas la caja existe sin abrir, que es un estado válido y visible.
      \item \textbf{\mbox{RN-10}} Una caja cobrada que no llega a crearse devuelve las monedas automáticamente.
    \end{cureglas}} \\ \hline
  \textbf{Observación} & \cuCelda{\textit{Sobre por qué la caja es una tabla y no un cobro directo.}\par
    Se consideró que comprar una caja sorteara y entregara en el mismo acto, sin dejar la caja como entidad. Se descartó porque el documento de diseño del juego prevé cajas ganadas jugando que caducan a los catorce días: eso exige que la caja exista sin abrir, con un origen y una fecha. Una vez que la caja tiene que existir, el estado \texttt{SIN\_\allowbreak{}ABRIR} sirve además para el \mbox{RN-10}, porque distingue el cobro realizado de la entrega pendiente, y para el \mbox{EX-05}, porque permite dejarla sin abrir cuando el sorteo falla en lugar de cobrarla y no dar nada.} \\ \hline
  \textbf{Datos que toca} & \cuCelda{\begin{cudatos}
      \item \textit{Sesion}: lee por token \textit{(paso 7)}
      \item \textit{TipoCaja}: lee precio, actividad y tabla de probabilidades \textit{(pasos 2, 8, 20)}
      \item \textit{TipoCajaObjeto}: lee los objetos posibles y su probabilidad \textit{(pasos 2, 20)}
      \item \textit{Caja}: crea, lee y marca como abierta o caducada \textit{(pasos 11, 18, 26)}
      \item \textit{CajaContenido}: crea tres \textit{(paso 24)}
      \item \textit{ObjetoCosmetico}: lee los candidatos del sorteo \textit{(paso 20)}
      \item \textit{UsuarioObjeto}: lee, para descartar repetidos; crea los nuevos \textit{(pasos 20, 22)}
      \item \textit{Usuario}: lee y escribe saldo y contador de cajas sin raro \textit{(pasos 9, 12, 23, 25)}
      \item \textit{MovimientoMoneda}: crea de tipo \texttt{COMPRA\_\allowbreak{}CAJA}, \texttt{CONVERSION} y, si procede, \texttt{DEVOLUCION} \textit{(pasos 12, 23)}
    \end{cudatos}} \\ \hline
\end{fichacu}
\clearpage
